\section*{Chapter \ref{ch:b1:cost-minimisation} --- Cost Minimisation}
% Numbers recomputed in code/bachelor-1/10-cost-minimisation/exercises/solutions.py (tested).

\begin{solution}{exo:b1:cost-minimisation:cd}
The tangency $K/L = w/r$ and $\sqrt{KL} = q$ give $L^* = q\sqrt{r/w}$ and
$K^* = q\sqrt{w/r}$. At $w = 4$, $r = 1$, $q = 10$: $L^* = 5$, $K^* = 20$, and the cost is
$4 \times 5 + 1 \times 20 = 40$ $(= 2q\sqrt{wr})$.
\end{solution}

\begin{solution}{exo:b1:cost-minimisation:leontief}
$K = 2 \times 100 = 200$ and $L = 3 \times 100 = 300$: $5 \times 200 + 10 \times 300 = 4000$. If $r$
doubles the inputs stay the same and the cost becomes $2000 + 3000 = 5000$.
\end{solution}

\begin{solution}{exo:b1:cost-minimisation:census}
$62/68 = 0.912$ and $10648/14971 = 0.711$ dalens of cost per dalen of output.
Not by itself: output is measured in value, not in units, so prices may
differ; the cost leaves out the cost of capital, which large firms use
more; and the two firms need not make the same goods. One would need firms
of different sizes making the same product, with all their costs.
\end{solution}

\begin{solution}{exo:b1:cost-minimisation:shift}
A rise in the wage shifts the \omterm{def:b1:cost-minimisation:cost-function}{cost function} ($C$ depends on $w$). A rise in
demand moves the firm along it, to a larger $q$. A new machine changes the
\omterm{def:b1:technology:production-function}{production function}, and so shifts the \omterm{def:b1:cost-minimisation:cost-function}{cost function} down.
\end{solution}

\begin{solution}{exo:b1:cost-minimisation:cd-cost}
The tangency gives $K/L = \tfrac{1/3}{2/3}\cdot\tfrac wr = \tfrac{w}{2r}$, and
$C(w, r, q) = q\,(3r)^{1/3}(\tfrac32w)^{2/3}$. At $w = 8$, $r = 1$:
$3^{1/3} \times 12^{2/3} = 7.5595$ per unit. Here $K/L = 4$, so $q = L \times 4^{1/3}$,
$L^* = 0.630q$ and $K^* = 2.520q$, and indeed $8 \times 0.630 + 2.520 = 7.56$.
Shephard: $\partial C/\partial w = \tfrac23C/w = \tfrac23 \times 7.5595/8 \times q = 0.630q = L^*$.
\end{solution}

\begin{solution}{exo:b1:cost-minimisation:short}
With $K = 4$, $L = q^2/4$ and $C_s = 4 + q^2/4$. In the long run $C = 2q\sqrt{wr} = 2q$.
At $q = 4$ both are 8 (four machine hours is the optimal capital for
$q = 4$); at $q = 6$, $C_s = 4 + 9 = 13 > 12$.
\end{solution}

\begin{solution}{exo:b1:cost-minimisation:mes-calc}
$AC = 100/q + 2 + 0.01q$ is lowest where $100/q^2 = 0.01$, at $q = 100$, with
$AC = 1 + 2 + 1 = 4$. $MC = 2 + 0.02q = 4$ there.
\end{solution}

\begin{solution}{exo:b1:cost-minimisation:three-doublings}
$0.8^3 = 0.512$: a fall of $48.8\%$.
\end{solution}

\begin{solution}{exo:b1:cost-minimisation:estimate}
From 50 to 400 is three doublings, over which cost was multiplied by
$120/200 = 0.6$. Per doubling, $0.6^{1/3} = 0.8434$: a \omterm{def:b1:cost-minimisation:learning-curve}{learning rate} of $15.7\%$.
\end{solution}

\begin{solution}{exo:b1:cost-minimisation:law}
By Shephard's lemma $L^* = \partial C/\partial w$, so
$\partial L^*/\partial w = \partial^2C/\partial w^2$. A concave function has nonpositive second
derivatives along each axis, so $\partial L^*/\partial w \le 0$. And
$\partial L^*/\partial r = \partial^2C/\partial r\,\partial w = \partial^2C/\partial w\,\partial r = \partial K^*/\partial w$,
since the order of differentiation does not matter.
\end{solution}

\begin{solution}{exo:b1:cost-minimisation:u-shape}
Each short-run curve lies on or above the long-run curve and touches it at
one output $q_0$ (\cref{prop:b1:cost-minimisation:envelope}), where the two
average costs have the same slope, because the total costs do and the
averages are both divided by $q_0$. Where long-run average cost falls, the
common slope is negative, so the short-run curve is still falling at $q_0$:
its minimum lies to the right. The envelope touches the minimum of a
short-run curve only at the bottom of the long-run U.
\end{solution}

\begin{solution}{exo:b1:cost-minimisation:separate}
One would need costs of many plants and firms over time, with their
cumulative output, their size and their research spending. \omterm{def:b1:cost-minimisation:learning-curve}{Learning by
doing} predicts that, at a given plant size and research spending, a firm's
cost falls with its own (or its industry's) cumulative output. Against it:
costs that fall at the same pace in years when cumulative output grows
slowly (a steady flow of research), or that are explained by plant size
alone, or new entrants with no experience whose costs match the
incumbents'.
\end{solution}

\begin{solution}{pb:b1:cost-minimisation:battery}
\textbf{1.} $f_L/f_K = K/L = w/r = 1/4$: the expansion path is $K = L/4$.

\textbf{2.} With $K = L/4$, $q = 2\sqrt{L^2/4} = L$: $L^* = q$ and $K^* = q/4$ (in general
$L^* = \tfrac q2\sqrt{r/w}$ and $K^* = \tfrac q2\sqrt{w/r}$). For $q = 1000$: 1,000 hours and 250
machine hours.

\textbf{3.} $C = q\sqrt{wr} = q\sqrt{10000} = 100q$: 100,000 dalens.

\textbf{4.} $AC = MC = 100$ dalens per kilowatt-hour: \omterm{def:b1:technology:returns-to-scale}{constant returns to scale}.

\textbf{5.} $\partial C/\partial w = \tfrac q2\sqrt{r/w} = \tfrac q2 \times 2 = q = L^*$.

\textbf{6.} $\sqrt{72 \times 200} = 120$ per unit; per unit,
$L^*/q = \tfrac12\sqrt{200/72} = 0.8333$ and $K^*/q = \tfrac12\sqrt{72/200} = 0.3$, so
$72 \times 0.8333 + 200 \times 0.3 = 120$. The plant uses less labour and more capital.

\textbf{7.} The cost doubles, to 200 per unit; the mix, which depends on $w/r$
only, does not change.

\textbf{8.} $L = q^2/(4 \times 250) = q^2/1000$, so $C_s = 200 \times 250 + 50q^2/1000 = 50000 + 0.05q^2$.

\textbf{9.} $SRAC = 50000/q + 0.05q$ and $SRMC = 0.1q$.

\textbf{10.} $50000/q^2 = 0.05$ at $q = 1000$, where $SRAC = 50 + 50 = 100$.

\textbf{11.} $SRAC = 41.67 + 60 = 101.67$, $SRMC = 120$, and the long-run average
cost is 100.

\textbf{12.} $62.5 + 40 = 102.5$.

\textbf{13.} In the long run the plant could choose any capital, including
250, so it can do no worse than with 250. They touch at $q = 1000$, the
output for which 250 is the cheapest capital ($K^* = q/4$).

\textbf{14.} 1,000; 2,000; 4,000; 8,000; 16,000.

\textbf{15.} $82.0/100 = 0.82$, $67.24/82.0 = 0.82$, $55.14/67.24 = 0.8200$,
$45.21/55.14 = 0.8199$. Cumulative output doubles each year, so the \omterm{def:b1:cost-minimisation:learning-curve}{learning
rate} is $18\%$.

\textbf{16.} $2^{-b} = 0.82$, $b = -\log_2 0.82 = 0.2863$.

\textbf{17.} From year 1 to year 5 cumulative output doubles four times:
$(45.21/100)^{1/4} = 0.820$, a \omterm{def:b1:cost-minimisation:learning-curve}{learning rate} of $18\%$.

\textbf{18.} Cumulative output doubled every year, so it and time moved
together: a fall of $18\%$ per doubling and a fall of $18\%$ a year predict the
same five numbers.

\textbf{19.} Unit cost is $2\sqrt{wr}/A = 200/A$, so $A = 200/45.21 = 4.42$.

\textbf{20.} $16000 + 3 \times 16000 = 64000$.

\textbf{21.} $64000/1000 = 64 = 2^6$: $100 \times 0.82^6 = 30.40$ dalens.

\textbf{22.} Seven years at $18\%$: $100 \times 0.82^7 = 24.93$ dalens.

\textbf{23.} If the cost falls through experience, the \omterm{def:b1:cost-minimisation:learning-curve}{learning curve}: output
stopped growing in year 5, so cumulative output doubles more slowly (only
twice from year 5 to year 8), and the cost falls more slowly. The yearly
rule would be right if the fall came from something steady in time, such
as research elsewhere. The five years cannot tell the two apart
(question 18); only years in which the growth of output changes, as from
year 5, can.

\textbf{24.} $1.1 \times 30.40 = 33.44$ dalens per kilowatt-hour.

\textbf{25.} The plant's \omterm{def:b1:cost-minimisation:learning-curve}{learning rate} is $18\%$ per doubling of cumulative
output, and its unit cost in year 8 is 30.40 dalens per kilowatt-hour.
\end{solution}
